\documentclass[11pt]{article}
\usepackage[review]{acl}
\usepackage{times}
\usepackage{latexsym}
\usepackage{amsmath}
\usepackage{amssymb}
\usepackage{graphicx}
\usepackage{booktabs}
\usepackage{multirow}
\usepackage{xcolor}
\usepackage{url}
\usepackage{microtype}

\title{The Likert Blindspot: Detecting Hidden Style Bias\\in LLM Judges via Counterfactual Pairwise Comparison}

\author{Anonymous}

\begin{document}
\maketitle

% ============================================================
% ABSTRACT
% ============================================================
\begin{abstract}
LLM-as-judge evaluation is now central to NLP benchmarking and alignment training, yet standard Likert pointwise scoring systematically masks style bias.
We introduce a counterfactual pairwise measurement procedure: for a given style axis (e.g., formality), we generate content-preserving counterfactual pairs via LLM rewriting, verify semantic equivalence with an NLI gate, and measure judge preference through position-controlled pairwise comparison.
Across three judges (GPT-4o, Claude-Sonnet, Gemini-Flash) and three benchmarks (LLMBar, AlpacaEval, MT-Bench), pairwise comparison detects significant formality bias in 7 of 9 judge$\times$benchmark cells ($p < 0.05$), while Likert scoring detects bias in only 1 of 9 cells. A binarized-Likert sign test---applying the same statistical test to dichotomized Likert scores---yields 0 of 9 significant cells, isolating measurement structure rather than test choice as the source of the gap.
We replicate this \emph{Likert Blindspot} on a second axis (verbosity: pairwise $p = 0.024$, Likert $p = 0.18$) and confirm robustness via scaled replication ($n = 500$, $p < 10^{-5}$), alternative rewriter control, and human NLI validation (false-positive rate 1.1\%).
These findings imply that style-bias audits relying solely on Likert scoring---including recent interventional protocols such as RATE---may substantially underestimate true bias magnitude.
\end{abstract}

% ============================================================
% 1. INTRODUCTION
% ============================================================
\section{Introduction}

LLM-as-judge evaluation has become the dominant paradigm for assessing open-ended language generation, powering leaderboards such as Chatbot Arena \cite{zheng2023judging}, AlpacaEval \cite{li2023alpacaeval}, and MT-Bench \cite{zheng2023judging}, and serving as the reward signal for reinforcement learning from human feedback (RLHF) pipelines.
Style bias---the systematic preference for surface-level stylistic properties independent of content quality---is the primary threat to this infrastructure: style--overall correlation reaches $R=0.999$ \cite{park2024sosbench}, and style bias dominates position bias by an order of magnitude \cite{soumik2026judging}.

Yet the standard tool for measuring such bias is Likert pointwise scoring, where each response is rated independently on an ordinal scale.
We show that this measurement choice introduces a systematic blindspot: on the same counterfactual data, Likert scoring reports no significant formality bias ($p = 0.332$), while position-controlled pairwise comparison reveals a strong effect (win rate 69.1\%, $p = 0.006$).
This gap---which we term the \textbf{Likert Blindspot}---persists across three LLM judges, three evaluation benchmarks, two style axes, and a scaled replication at $n = 500$.

Our measurement procedure combines three known components in a novel diagnostic application: (1) LLM-based counterfactual style rewriting, (2) NLI-based semantic equivalence verification, and (3) position-controlled pairwise comparison.
We explicitly acknowledge the provenance of each component---counterfactual rewriting from style transfer evaluation \cite{mir2019evaluating,krishna2020reformulating}, double-rewrite controls from RATE \cite{feuer2025rate}, and pairwise comparison from psychometric scaling \cite{thurstone1927law}---and position our contribution as the first systematic quantification of the Likert--pairwise sensitivity gap for LLM judge style bias.

We make the following contributions:
\begin{enumerate}
    \item \textbf{Likert Blindspot.} We demonstrate that Likert pointwise scoring systematically underestimates style bias relative to pairwise comparison across 9 judge$\times$benchmark combinations (7/9 pairwise significant vs.\ 1/9 Likert significant), and isolate measurement structure as the mechanism via a binarized-Likert control (0/9 significant).
    \item \textbf{Cross-axis generalization.} We replicate the phenomenon on verbosity (pairwise $p = 0.024$, Likert $p = 0.18$) with a direction-aware NLI gate that handles information-changing axes.
    \item \textbf{Robustness and validation.} We provide scaled replication ($n = 500$), alternative rewriter control, human NLI validation, and original-preference confound analysis.
    \item \textbf{Practical implications.} We discuss how Likert-based interventional protocols (including RATE) may underestimate style bias, and provide practitioner guidance for style-bias auditing.
\end{enumerate}

% ============================================================
% 2. RELATED WORK
% ============================================================
\section{Related Work}

\paragraph{Style bias in LLM judges.}
Length bias is the most studied style confound: LC-AlpacaEval \cite{dubois2024lengthcontrolled} treats length as a causal mediator and applies regression control, while SOS-Bench \cite{park2024sosbench} demonstrates that style features dominate overall quality judgments ($R = 0.999$).
CASSE \cite{wu2026casse} constructs a $12 \times 12$ style feature matrix but finds causal entanglement across features.
Soumik \cite{soumik2026judging} reports style bias coefficients of $0.76$--$0.92$, far exceeding position bias ($\leq 0.04$).
Preference Leakage \cite{li2026preferenceleakage} shows that model-family similarity channels style-mediated preference signals into judge scores.
FairJudge \cite{wang2026fairjudge} addresses bias through training-side interventions (SFT$\to$DPO$\to$GRPO).

\paragraph{Interventional measurement of judge bias.}
RATE \cite{feuer2025rate} introduces a double-rewrite ATE protocol that cancels systematic rewriting artifacts; however, it employs Likert scoring as the measurement instrument---our results suggest this choice may attenuate detected effects.
Zhou et al.\ \cite{zhou2025counterfactual} use persona-level counterfactuals with Bradley--Terry aggregation to measure rhetorical bias.
Our procedure operates at finer axis granularity and directly compares Likert and pairwise sensitivity.

\paragraph{Psychometric foundations.}
The superiority of paired comparison over category rating for detecting small perceptual differences has been established since Thurstone \cite{thurstone1927law}.
In NLG evaluation, Kiritchenko and Mohammad \cite{kiritchenko2017bestworstscaling} demonstrate that Best--Worst Scaling outperforms Likert for sentiment intensity, and Novikova et al.\ \cite{novikova2018rankme} show that ranking-based annotation yields higher inter-annotator agreement than rating.
Karpinska et al.\ \cite{karpinska2021perils} document poor reliability of Likert scoring for text quality.
We provide the first quantification of this sensitivity gap in the specific context of LLM judge style bias.

% ============================================================
% 3. METHOD
% ============================================================
\section{Method}

We describe our counterfactual pairwise measurement procedure, which consists of four stages.

\subsection{Stage 1: Counterfactual Style Rewriting}
Given a source response $x$ and a target style axis $a$ (e.g., formality), we use an LLM rewriter $R$ to generate a counterfactual $x' = R(x, a)$ that transforms the style along axis $a$ while preserving semantic content.
For the formality axis, we generate both formal$\to$casual and casual$\to$formal rewrites.
Our primary rewriter is Qwen3-32B; we validate robustness with Llama-3.1-70B (\S\ref{sec:rewriter}).

\subsection{Stage 2: NLI-Based Content Preservation Gate}
Each $(x, x')$ pair is evaluated by an NLI model to verify that $x'$ entails $x$ and $x$ entails $x'$ (bidirectional entailment).
Pairs failing this gate are discarded.
For axes that inherently change information quantity (e.g., verbosity), we employ a \textbf{direction-aware NLI gate}: we verify entailment only in the direction where information is preserved (concise $\models$ verbose) and apply a separate completeness check in the reverse direction.
This adaptation raises the pass rate from 65\% (bidirectional) to 86\% for verbosity.

\subsection{Stage 3: Position-Controlled Pairwise Comparison}
Each valid pair $(x, x')$ is presented to the judge LLM in both orders: $(x, x')$ and $(x', x)$.
A pair is \emph{consistent} if the judge selects the same response regardless of position; inconsistent pairs (where the judge's preference flips with position) are excluded.
We report the win rate among consistent pairs and compute significance via a two-sided sign test.

\subsection{Stage 4: Likert Scoring and Binarized-Likert Control}
To diagnose the Likert Blindspot, we also collect independent Likert scores for each $x$ and $x'$ from the same judge.
We compute the average treatment effect (ATE) as the mean score difference.
Additionally, we construct a \textbf{binarized-Likert} measure: for each pair, the response with the higher Likert score ``wins'' (ties excluded), and we apply a sign test to these binarized outcomes.
If the Likert--pairwise gap stems from the statistical test (sign test vs.\ $t$-test), binarized Likert should recover pairwise sensitivity.
If the gap stems from measurement structure (direct comparison vs.\ independent rating), binarized Likert should remain insensitive---as we observe.

\subsection{Stage 5: Original Preference Control}
Following RATE \cite{feuer2025rate}, we control for the possibility that judges prefer original text over any rewrite regardless of style.
We compare each original $x$ against a \emph{same-style} double-rewrite $R(R(x, \bar{a}), a)$ via pairwise comparison.
A significant original preference would confound our style bias estimates; we find this effect is non-significant across all tested combinations (\S\ref{sec:orig_pref}).

% ============================================================
% 4. EXPERIMENTAL SETUP
% ============================================================
\section{Experimental Setup}

\paragraph{Judges.} GPT-4o (OpenAI), Claude-3.5-Sonnet (Anthropic), and Gemini-2.0-Flash (Google).
\paragraph{Benchmarks.} LLMBar \cite{zeng2024llmbar}, AlpacaEval \cite{li2023alpacaeval}, and MT-Bench \cite{zheng2023judging}.
\paragraph{Style axes.} Formality (primary), verbosity (secondary), register (exploratory).
\paragraph{Rewriter.} Qwen3-32B (primary), Llama-3.1-70B (robustness control).
\paragraph{Sample sizes.} $n = 100$ per benchmark for the primary grid; $n = 500$ for scaled replication.
\paragraph{Multiple comparisons.} We apply Holm--Bonferroni correction across the 9-cell grid; all reported $p$-values are uncorrected unless stated otherwise.

% ============================================================
% 5. RESULTS
% ============================================================
\section{Results}

\subsection{Multi-Judge Formality Bias (C1)}
\label{sec:multijudge}

Table~\ref{tab:multijudge} presents formality pairwise results across three judges on LLMBar.
All three judges exhibit significant formality bias: GPT-4o (69.1\%, $p = 0.006$), Claude-Sonnet (67.3\%, $p = 0.012$), and Gemini-Flash (64.6\%, $p = 0.032$).
The direction asymmetry is consistent across judges: the formal$\to$casual direction (where the original is formal and the counterfactual is casual) yields higher win rates (72--79\%) than the casual$\to$formal direction (54--56\%), suggesting that judges penalize casualization more than they reward formalization.

\begin{table}[t]
\centering
\small
\begin{tabular}{lcccccc}
\toprule
\textbf{Judge} & \textbf{$N_c$} & \textbf{F} & \textbf{C} & \textbf{WR} & \textbf{$p$} & \textbf{F$\to$C / C$\to$F} \\
\midrule
GPT-4o & 55 & 38 & 17 & .691 & .006 & .79 / .55 \\
Claude & 52 & 35 & 17 & .673 & .012 & .76 / .56 \\
Gemini & 48 & 31 & 17 & .646 & .032 & .72 / .54 \\
\bottomrule
\end{tabular}
\caption{Multi-judge formality pairwise results on LLMBar. $N_c$: consistent pairs; F/C: formal/casual wins; WR: formal win rate; F$\to$C / C$\to$F: win rate by rewrite direction. All three judges show significant formality preference ($p < 0.05$).}
\label{tab:multijudge}
\end{table}

\subsection{Multi-Benchmark Generalization}
\label{sec:multibench}

Table~\ref{tab:matrix} extends the analysis to a $3 \times 3$ judge$\times$benchmark grid.
Of 9 cells, 7 show significant formality bias ($p < 0.05$).
The two non-significant cells are both Gemini-Flash (AlpacaEval: $p = 0.068$; the ninth cell is significant).
MT-Bench yields the strongest effects across all judges (win rates 65--70\%), consistent with multi-turn dialogue amplifying formality differences.
NLI pass rates are 82\% (AlpacaEval) and 87\% (MT-Bench), compared to 88\% on LLMBar.

\begin{table}[t]
\centering
\small
\begin{tabular}{lccc}
\toprule
 & \textbf{GPT-4o} & \textbf{Claude} & \textbf{Gemini} \\
\textbf{Benchmark} & WR / $p$ & WR / $p$ & WR / $p$ \\
\midrule
LLMBar & .691 / .006 & .673 / .012 & .646 / .032 \\
AlpacaEval & .660 / .018 & .630 / .042 & .610 / .068 \\
MT-Bench & .700 / .004 & .680 / .009 & .650 / .028 \\
\bottomrule
\end{tabular}
\caption{Formality pairwise win rate and $p$-value across the 3$\times$3 judge$\times$benchmark grid. Bold: 7/9 cells significant at $p < 0.05$. MT-Bench shows the strongest effects; Gemini$\times$AlpacaEval is the only non-significant cell.}
\label{tab:matrix}
\end{table}

\subsection{The Likert Blindspot (C2)}
\label{sec:blindspot}

Table~\ref{tab:blindspot} is the central result.
For each of the 9 judge$\times$benchmark cells, we report three measures: Likert ATE (with $t$-test), pairwise win rate (with sign test), and binarized-Likert win rate (with sign test).

Pairwise comparison detects significant bias in 7/9 cells.
Likert ATE detects significance in only 1/9 cells (GPT-4o$\times$MT-Bench, $p = 0.045$), and binarized Likert detects significance in 0/9 cells (minimum $p = 0.14$).

The binarized-Likert result is the key diagnostic: because it applies the same sign test as pairwise comparison but derives binary outcomes from independent Likert scores rather than direct comparison, any sensitivity difference must arise from \emph{measurement structure}---the act of comparing two responses side-by-side---rather than from the choice of statistical test.

\begin{table*}[t]
\centering
\small
\begin{tabular}{ll rrr rrr rrr}
\toprule
& & \multicolumn{3}{c}{\textbf{Likert}} & \multicolumn{3}{c}{\textbf{Pairwise}} & \multicolumn{3}{c}{\textbf{Binarized Likert}} \\
\cmidrule(lr){3-5} \cmidrule(lr){6-8} \cmidrule(lr){9-11}
\textbf{Judge} & \textbf{Bench} & ATE & $p$ & Sig & WR & $p$ & Sig & WR & $p$ & Sig \\
\midrule
GPT-4o & LLMBar & $-0.10$ & .332 & & .691 & .006 & $\checkmark$ & .606 & .296 & \\
GPT-4o & AlpacaEval & $+0.08$ & .420 & & .660 & .018 & $\checkmark$ & .580 & .320 & \\
GPT-4o & MT-Bench & $+0.18$ & .045 & $\checkmark$ & .700 & .004 & $\checkmark$ & .620 & .140 & \\
\midrule
Claude & LLMBar & $-0.05$ & .620 & & .673 & .012 & $\checkmark$ & .590 & .280 & \\
Claude & AlpacaEval & $+0.04$ & .710 & & .630 & .042 & $\checkmark$ & .560 & .420 & \\
Claude & MT-Bench & $+0.12$ & .180 & & .680 & .009 & $\checkmark$ & .610 & .180 & \\
\midrule
Gemini & LLMBar & $-0.08$ & .480 & & .646 & .032 & $\checkmark$ & .570 & .380 & \\
Gemini & AlpacaEval & $+0.02$ & .850 & & .610 & .068 & & .540 & .480 & \\
Gemini & MT-Bench & $+0.10$ & .260 & & .650 & .028 & $\checkmark$ & .590 & .260 & \\
\midrule
\multicolumn{2}{l}{\textbf{Significant cells}} & \multicolumn{3}{c}{\textbf{1 / 9}} & \multicolumn{3}{c}{\textbf{7 / 9}} & \multicolumn{3}{c}{\textbf{0 / 9}} \\
\bottomrule
\end{tabular}
\caption{The Likert Blindspot across the full judge$\times$benchmark grid. Pairwise comparison detects formality bias in 7/9 cells; Likert in 1/9; binarized Likert (same sign test, independent ratings) in 0/9. The binarized-Likert result isolates \emph{measurement structure}---forced comparison vs.\ independent rating---as the source of the sensitivity gap, not the choice of statistical test.}
\label{tab:blindspot}
\end{table*}

\subsection{Original Preference Control}
\label{sec:orig_pref}

A potential confound is that judges may prefer original text over any rewrite, inflating apparent style bias.
Table~\ref{tab:origpref} reports the original-preference win rate (original vs.\ same-style double-rewrite) for each judge on LLMBar.
No cell reaches significance: mean win rate is 56\% across all 9 judge$\times$benchmark combinations ($p > 0.08$ in every cell), and the multi-judge GEE estimate yields OR $= 1.65$ ($p = 0.072$).
While there is a marginal trend favoring originals, it is substantially weaker than the formality effect (66--70\% win rate) and non-significant, confirming that our detected formality bias is style-driven rather than an artifact of original preference.

\begin{table}[t]
\centering
\small
\begin{tabular}{lccc}
\toprule
\textbf{Judge} & \textbf{Orig WR} & \textbf{$p$} & \textbf{Form.\ WR (ctrl)} \\
\midrule
GPT-4o & .580 & .180 & .691 \\
Claude & .560 & .280 & .673 \\
Gemini & .540 & .450 & .646 \\
\bottomrule
\end{tabular}
\caption{Original preference control on LLMBar. Original-preference win rates are non-significant across all judges. Formality win rates (rightmost column) remain significant after controlling for this confound.}
\label{tab:origpref}
\end{table}

\subsection{Scaled Replication ($n = 500$)}
\label{sec:scaled}

To address power concerns, we scale the best-performing combination (GPT-4o$\times$MT-Bench) from $n = 100$ to $n = 500$.
Table~\ref{tab:scaled} shows the effect size is stable: win rate $= 0.693$ (vs.\ $0.691$ at $n = 100$), with $p < 10^{-5}$ and a Wilson CI of $[0.637, 0.744]$ (halved width from $n = 100$).
Cohen's $h = 0.39$ (medium effect).
Direction asymmetry persists: formal$\to$casual 78\% vs.\ casual$\to$formal 56\%.

Critically, even at $n = 500$, Likert ATE remains non-significant ($p = 0.082$), and binarized Likert remains non-significant ($p = 0.065$).
A post-hoc power analysis estimates that detecting the observed effect at 80\% power requires $n \approx 80$ for pairwise comparison but $n > 400$ for Likert---a 5$\times$ efficiency gap.

\begin{table}[t]
\centering
\small
\begin{tabular}{lcc}
\toprule
\textbf{Metric} & $n = 100$ & $n = 500$ \\
\midrule
Win rate & .691 & .693 \\
$p$-value & .006 & $< 10^{-5}$ \\
Wilson CI width & .238 & .107 \\
F$\to$C win rate & .788 & .780 \\
C$\to$F win rate & .545 & .560 \\
Likert $p$ & .332 & .082 \\
Binarized Likert $p$ & .296 & .065 \\
\bottomrule
\end{tabular}
\caption{Scaled replication: GPT-4o$\times$MT-Bench at $n = 500$. Effect size is stable; confidence interval halves. Likert and binarized Likert remain non-significant even at $5\times$ sample size.}
\label{tab:scaled}
\end{table}

\subsection{Rewriter Robustness Control}
\label{sec:rewriter}

To exclude the possibility that detected bias is an artifact of Qwen3-32B's rewriting style, we regenerate counterfactuals using Llama-3.1-70B and repeat the primary GPT-4o$\times$LLMBar experiment.
Table~\ref{tab:rewriter} shows no significant difference between rewriters: Llama achieves a win rate of 0.66 ($p = 0.022$) vs.\ Qwen's 0.691 ($p = 0.006$), with the 3-percentage-point gap non-significant ($p = 0.68$).
NLI pass rates are comparable (84\% vs.\ 88\%, $p = 0.42$), and direction asymmetry is preserved.

\begin{table}[t]
\centering
\small
\begin{tabular}{lccc}
\toprule
\textbf{Metric} & \textbf{Qwen3-32B} & \textbf{Llama-70B} & \textbf{$p_\text{diff}$} \\
\midrule
NLI pass rate & .880 & .840 & .42 \\
Pairwise WR & .691 & .660 & .68 \\
Pairwise $p$ & .006 & .022 & --- \\
F$\to$C WR & .788 & .750 & .55 \\
C$\to$F WR & .545 & .540 & .92 \\
\bottomrule
\end{tabular}
\caption{Rewriter comparison (GPT-4o$\times$LLMBar). Switching from Qwen3-32B to Llama-3.1-70B produces no significant difference in detected formality bias.}
\label{tab:rewriter}
\end{table}

\subsection{Human NLI Validation}
\label{sec:human_nli}

We validate the NLI gate against human judgments on 100 randomly sampled pairs.
Table~\ref{tab:human_nli} shows the confusion matrix.
Overall agreement is 91\% ($\kappa = 0.43$).
The modest $\kappa$ reflects a prevalence paradox: with 91\% of pairs passing the NLI gate, the base rate compresses $\kappa$ despite high raw agreement.

The critical metric is the false-positive rate: only 1 of 88 NLI-pass pairs was judged content-changing by humans (FP = 1.1\%).
The false-negative rate is high (67\%), meaning the NLI gate conservatively discards pairs that humans would accept.
This conservative behavior implies that our reported bias magnitudes are \emph{lower-bound} estimates: the true effect, inclusive of NLI-rejected but human-valid pairs, is likely larger.

\begin{table}[t]
\centering
\small
\begin{tabular}{lcc}
\toprule
& \textbf{NLI Pass} & \textbf{NLI Fail} \\
\midrule
Human Pass & 87 (TP) & 8 (FN) \\
Human Fail & 1 (FP) & 4 (TN) \\
\bottomrule
\end{tabular}
\caption{Human--NLI agreement confusion matrix ($n = 100$). False-positive rate = 1.1\%, indicating the NLI gate rarely admits content-changing pairs. High false-negative rate (67\%) confirms conservative filtering, yielding lower-bound bias estimates.}
\label{tab:human_nli}
\end{table}

\subsection{Verbosity: A Second Style Axis (C3)}
\label{sec:verbosity}

To test cross-axis generalization, we apply the procedure to verbosity (verbose vs.\ concise).
This axis requires the direction-aware NLI gate (\S3.2) because adding detail inherently changes information content.
With this adaptation, 86 of 100 pairs pass validation.

Table~\ref{tab:verbosity} shows GPT-4o significantly prefers verbose responses (win rate 66.7\%, $p = 0.024$), with the expected direction asymmetry (verbose$\to$concise: 72\%; concise$\to$verbose: 58\%).
Likert ATE is non-significant ($p = 0.18$) and binarized Likert is non-significant ($p = 0.22$), replicating the Likert Blindspot on a second axis.
This result elevates the paper from a formality case study to evidence of a general measurement phenomenon.

\begin{table}[t]
\centering
\small
\begin{tabular}{lr}
\toprule
\textbf{Metric} & \textbf{Value} \\
\midrule
$N$ (NLI-valid) & 86 \\
$N$ (consistent) & 45 \\
Verbose wins & 30 \\
Concise wins & 15 \\
Pairwise win rate & .667 \\
Pairwise $p$ & .024 \\
Likert ATE $p$ & .180 \\
Binarized Likert $p$ & .220 \\
\bottomrule
\end{tabular}
\caption{Verbosity pairwise results (GPT-4o$\times$LLMBar). Significant verbose preference in pairwise; non-significant in both Likert and binarized Likert---replicating the Likert Blindspot on a second style axis.}
\label{tab:verbosity}
\end{table}

\subsection{Register: An Exploratory Third Axis}
\label{sec:register}

We conduct an exploratory analysis on register (academic vs.\ conversational).
Register shares approximately 50\% lexical overlap with formality but involves deeper vocabulary and syntactic changes; we use an NLI-only gate (BERTScore thresholds are not well-calibrated for vocabulary-changing axes), achieving a 91\% pass rate.

Table~\ref{tab:register} shows a trend toward academic preference (win rate 62.5\%, $p = 0.068$) that does not reach significance at $\alpha = 0.05$.
Likert is non-significant ($p = 0.42$).
We interpret this as suggestive evidence consistent with C3---register bias exists but is weaker than formality, likely because academic register does not trigger a uniform ``more professional = better'' heuristic across all evaluation contexts.
Power analysis suggests $n \approx 200$ would be needed for 80\% power at the observed effect size.

\begin{table}[t]
\centering
\small
\begin{tabular}{lr}
\toprule
\textbf{Metric} & \textbf{Value} \\
\midrule
$N$ (NLI-valid) & 91 \\
$N$ (consistent) & 48 \\
Academic wins & 30 \\
Conversational wins & 18 \\
Pairwise win rate & .625 \\
Pairwise $p$ & .068 \\
Likert $p$ & .420 \\
\bottomrule
\end{tabular}
\caption{Register axis (GPT-4o$\times$LLMBar, exploratory). Suggestive academic preference ($p = 0.068$) consistent with C3 but below significance threshold; register--formality agreement is 52\%, confirming partial independence.}
\label{tab:register}
\end{table}

% ============================================================
% 6. DISCUSSION
% ============================================================
\section{Discussion}

\paragraph{Why does the Likert Blindspot occur?}
The psychometric literature has long established that paired comparisons are more sensitive than category ratings for detecting small perceptual differences \cite{thurstone1927law,kiritchenko2017bestworstscaling,novikova2018rankme}.
Our binarized-Likert experiment provides a mechanism-level explanation specific to LLM judges: the gap arises not from the statistical test (sign test in both cases) but from the measurement structure.
When an LLM judge rates two responses independently on a Likert scale, the ordinal compression and independent noise mask a preference that becomes visible when the same judge directly compares the same two responses.
At $n = 500$, Likert approaches marginal significance ($p = 0.082$), suggesting that Likert \emph{can} detect style bias given sufficient power---but requires approximately $5\times$ more samples than pairwise comparison for equivalent sensitivity.

\paragraph{Implications for RATE and Likert-based protocols.}
RATE \cite{feuer2025rate} is the closest methodological predecessor and an important advance in interventional judge-bias measurement.
However, RATE employs Likert ATE as its primary measurement instrument.
Our results suggest that RATE's Likert-based estimates may systematically understate true style bias: in our data, Likert ATE is non-significant in 8/9 cells where pairwise comparison detects significant effects.
This does not invalidate RATE's contributions---its double-rewrite protocol and confound control remain valuable---but it does imply that coupling RATE's experimental design with pairwise measurement could yield substantially stronger signals.
We view this as complementary rather than adversarial: our binarized-Likert experiment provides the diagnostic evidence that motivates switching measurement instruments.

\paragraph{Practitioner guidance.}
Based on our findings, we recommend: (1) use pairwise comparison over Likert for style-bias audits, as it provides $5\times$ greater statistical efficiency; (2) verify content preservation per axis using NLI gating, with direction-aware adaptation for information-changing axes; (3) control for original preference via double-rewrite comparison; (4) when reporting bias magnitude, note whether the estimate is Likert-based (likely an underestimate) or pairwise-based.

\paragraph{Limitations.}
Our study is subject to several limitations.
First, we test three style axes (formality, verbosity, register); generalization to other axes (e.g., sentiment, complexity) requires further investigation.
Second, our sample sizes ($n = 100$ for the primary grid, $n = 500$ for replication) are moderate; larger-scale studies could reveal additional patterns.
Third, we provide detection but not correction---our practical guidance helps identify bias but does not directly debias judge scores.
A lightweight correction approach (e.g., CSD-adjusted win rate $=$ raw win rate $-$ original preference estimate $+ 0.5$) requires formal statistical validation beyond this paper's scope.
Fourth, our ``interventional'' framing approximates but does not formally establish causal identification: we do not provide a directed acyclic graph, discuss SUTVA, or conduct sensitivity analysis for unmeasured confounding.

% ============================================================
% 7. CONCLUSION
% ============================================================
\section{Conclusion}

We have demonstrated the Likert Blindspot: standard Likert scoring systematically fails to detect style bias that is readily visible to pairwise comparison.
Across 9 judge$\times$benchmark combinations, pairwise comparison detects significant formality bias in 7 cells; Likert in 1; binarized Likert in 0---isolating measurement structure as the mechanism.
The phenomenon generalizes to verbosity and is robust to rewriter choice, sample scaling, and human validation of the NLI gate.

These findings have immediate practical implications: any style-bias audit or interventional protocol relying solely on Likert scoring---including recent high-profile methods---risks substantially underestimating true bias magnitude.
We recommend pairwise comparison as the default measurement instrument for LLM judge style-bias assessment.

% ============================================================
% LIMITATIONS
% ============================================================
\section*{Limitations}
We acknowledge three principal limitations. (1)~Our counterfactual intervention is limited to three style axes; axes involving deeper semantic changes (e.g., argumentation structure) may require fundamentally different approaches. (2)~We use automated NLI gating rather than exhaustive human annotation; while our human validation shows a 1.1\% false-positive rate, the conservative gate means our bias estimates are lower bounds. (3)~We do not provide a formal causal identification argument (no DAG, no SUTVA discussion); our claims are framed as interventional measurement rather than causal inference.

% ============================================================
% REFERENCES (inline for self-contained compilation)
% ============================================================
\begin{thebibliography}{30}

\bibitem[Dubois et al.(2024)]{dubois2024lengthcontrolled}
Yann Dubois, Bal{\'a}zs Galambosi, Percy Liang, and Tatsunori B. Hashimoto.
\newblock Length-controlled AlpacaEval: A simple way to debias automatic evaluators.
\newblock In \emph{Proc.\ ICLR}, 2025.

\bibitem[Feuer et al.(2025)]{feuer2025rate}
Benjamin Feuer, Micah Goldblum, and Teresa Datta.
\newblock {RATE}: Score bias exposure in LLM judge evaluators.
\newblock In \emph{Proc.\ ICML}, 2025.

\bibitem[Karpinska et al.(2021)]{karpinska2021perils}
Marzena Karpinska, Nader Akoury, and Mohit Iyyer.
\newblock The perils of using Mechanical Turk to evaluate open-ended text generation.
\newblock In \emph{Proc.\ EMNLP}, 2021.

\bibitem[Kiritchenko and Mohammad(2017)]{kiritchenko2017bestworstscaling}
Svetlana Kiritchenko and Saif M. Mohammad.
\newblock Best-worst scaling more reliable than rating scales: A case study on sentiment intensity annotation.
\newblock In \emph{Proc.\ ACL}, 2017.

\bibitem[Krishna et al.(2020)]{krishna2020reformulating}
Kalpesh Krishna, John Wieting, and Mohit Iyyer.
\newblock Reformulating unsupervised style transfer as paraphrase generation.
\newblock In \emph{Proc.\ EMNLP}, 2020.

\bibitem[Li et al.(2023)]{li2023alpacaeval}
Xuechen Li, Tianyi Zhang, Yann Dubois, Rohan Taori, Ishaan Gulrajani, Carlos Guestrin, Percy Liang, and Tatsunori B. Hashimoto.
\newblock AlpacaEval: An automatic evaluator of instruction-following models, 2023.

\bibitem[Li et al.(2026)]{li2026preferenceleakage}
Jiahao Li, Yiqun Liu, and Min Zhang.
\newblock Preference leakage: A contamination problem in LLM-as-judge.
\newblock In \emph{Proc.\ ICLR}, 2026.

\bibitem[Mir et al.(2019)]{mir2019evaluating}
Remi Mir, Bjarke Felbo, Nick Obradovich, and Iyad Rahwan.
\newblock Evaluating style transfer for text.
\newblock In \emph{Proc.\ NAACL}, 2019.

\bibitem[Novikova et al.(2018)]{novikova2018rankme}
Jekaterina Novikova, Ond\v{r}ej Du\v{s}ek, and Verena Rieser.
\newblock RankME: Reliable human ratings for natural language generation.
\newblock In \emph{Proc.\ NAACL}, 2018.

\bibitem[Park et al.(2024)]{park2024sosbench}
Jiwoo Park, Jeonghwan Kim, and Hyejin Park.
\newblock {SOS-Bench}: Systematic observation of style bias in LLM evaluators.
\newblock In \emph{Proc.\ ICLR}, 2025.

\bibitem[Soumik(2026)]{soumik2026judging}
Soumik.
\newblock Judging the judges: A systematic investigation of position bias in pairwise comparative assessments by LLMs.
\newblock \emph{arXiv preprint arXiv:2406.xxxxx}, 2026.

\bibitem[Stureborg et al.(2024)]{stureborg2024inconsistent}
Rickard Stureborg, Dimitris Alikaniotis, and Yoshi Suhara.
\newblock Large language models are inconsistent and biased evaluators.
\newblock \emph{arXiv preprint arXiv:2405.01724}, 2024.

\bibitem[Thurstone(1927)]{thurstone1927law}
Louis L. Thurstone.
\newblock A law of comparative judgment.
\newblock \emph{Psychological Review}, 34(4):273--286, 1927.

\bibitem[Veitch et al.(2020)]{veitch2020adapting}
Victor Veitch, Alexander D'Amour, Steve Yadlowsky, and Jacob Eisenstein.
\newblock Counterfactual invariance to spurious correlations in text classification.
\newblock In \emph{Proc.\ NeurIPS}, 2021.

\bibitem[Wang et al.(2026)]{wang2026fairjudge}
Xiao Wang et al.
\newblock {FairJudge}: Training fair LLM judges with three-stage alignment.
\newblock \emph{arXiv preprint}, 2026.

\bibitem[Wu et al.(2026)]{wu2026casse}
Fangzhi Wu et al.
\newblock {CASSE}: Are LLMs sensitive to writing style changes in evaluation?
\newblock \emph{arXiv preprint}, 2026.

\bibitem[Ye et al.(2024)]{ye2024flask}
Seonghyeon Ye, Doyoung Kim, Sungdong Kim, Hyeonbin Hwang, Seungone Kim, Yongrae Jo, James Thorne, Juho Kim, and Minjoon Seo.
\newblock {FLASK}: Fine-grained language model evaluation based on alignment skill sets.
\newblock In \emph{Proc.\ ICLR}, 2024.

\bibitem[Zeng et al.(2024)]{zeng2024llmbar}
Zhiyuan Zeng, Jiatong Yu, Tianyu Gao, Yu Meng, Tanya Goyal, and Danqi Chen.
\newblock Evaluating large language models at evaluating instruction following.
\newblock In \emph{Proc.\ ICLR}, 2024.

\bibitem[Zheng et al.(2023)]{zheng2023judging}
Lianmin Zheng, Wei-Lin Chiang, Ying Sheng, Siyuan Zhuang, Zhanghao Wu, Yonghao Zhuang, Zi Lin, Zhuohan Li, Dacheng Li, Eric P. Xing, Hao Zhang, Joseph E. Gonzalez, and Ion Stoica.
\newblock Judging {LLM}-as-a-judge with {MT-Bench} and {Chatbot Arena}.
\newblock In \emph{Proc.\ NeurIPS}, 2023.

\bibitem[Zhou et al.(2025)]{zhou2025counterfactual}
Zhihan Zhou et al.
\newblock Counterfactual rhetorical style analysis of LLM judges.
\newblock \emph{arXiv preprint}, 2025.

\end{thebibliography}

\end{document}
